\documentclass[10pt,a4paper]{article}
\usepackage{fontspec,amsmath,amssymb,graphicx}
\usepackage[margin=18mm]{geometry}
\setmainfont{Times New Roman}
\setlength{\parindent}{0pt}
\setlength{\parskip}{5pt}
\allowdisplaybreaks
\sloppy
\emergencystretch=2em
\begin{document}
\section*{BT\_6\_1\_1\_2\_05 — Lời giải bổ sung}
Tích thành tổng cho tích phân sine bằng
\[\frac{\sin((n-m)\pi/2)}{n-m}-\frac{\sin((n+m)\pi/2)}{n+m},\]
và cosine thay dấu trừ thành cộng. Vì $n-m$ chẵn nên $n+m$ cũng chẵn, hai tử bằng 0. Tuy nhiên đề cho $m,n\in\mathbb Z$ và chỉ loại $n=m$, chưa loại $n=-m$. Nếu $n=-m\ne0$, hai hàm cùng tần số và tích phân sine bằng $-\pi/2$, cosine bằng $\pi/2$, không bằng 0. Kết luận đúng với $n\ne\pm m$ (hoặc $m,n$ không âm khác nhau). Trường hợp chỉ số 0 xử lý trực tiếp; sine 0 đồng nhất 0.\end{document}